\section{Donna, nel mio ritorno}

\textbf{Torquato Tasso}\footnote{\url{https://en.wikipedia.org/wiki/Jerusalem\_Delivered}} was a Renaissance Italian poet well known for his love poems, the epic poem \emph{Jerusalem Delivered}\footnote{\url{https://en.wikipedia.org/wiki/Jerusalem\_Delivered}} about the Crusades, and a blank verse retelling of the seven days of creation in Genesis. His work brought him fame in his day and, until recently, he was one of the most read poets in Europe. Many subsequent poets were inspired by him and several madrigals were composed for his poems. This, composed by \textbf{Claudio Monteverdi}, is one of them.

Unfortunately, Tasso drifted into insanity. Giacomo Leopardi wrote Dialogo di Torquato Tasso e del suo Genio familiare\footnote{\url{https://www.gornahoor.net/?p=25}} about his stay at the St. Anna insane asylum. This is a dialogue between Tasso and his family spirit.

\columnratio{0.4}
\begin{paracol}{2}
\begin{verse}
Donna, nel mio ritorno\\
il mio pensiero, a cui nulla pon freno,\\
precorre dove il ciel è più sereno,\\
e se ne vien a far con voi soggiorno;\\
né da voi si diparte\\
giamai la notte e'l giorno,\\
perché l'annoia ciascun'altra parte.\\
Onde sol per virtù del pensier mio,\\
mentre ne vengo a voi, con voi son io.
\end{verse}
\switchcolumn
\begin{verse}
Milady, as I return to you\\
my thoughts, which nothing can restrain,\\
rush across the cloudless sky\\
and come to abide with you;\\
never do they leave you,\\
neither by day nor night,\\
because any other thought would bore me.\\
Thus, thanks only to my thoughts\\
As I approach you, I am already with you.
\end{verse}
\end{paracol}
\flright{\textsc{Torquato Tasso}, 1590}

This is my own translation, based on personal experience.

\url{https://www.youtube.com/watch?v=oxJ2Zv8PGbo} 

\flrightit{Posted on 2021-12-27 by Cologero}

\begin{center}* * *\end{center}

\begin{footnotesize}
\begin{sffamily}
\texttt{Arthur Konrad on 2021-12-28 at 04:25 said:}

Venetian music (and especially opera) is second to none. Baroque has experienced quite a bit of revival, in every sense, from performances to instrument construction. Venetian libretti are always splendid -- how different an air from the latter day defilement of poetry with intellectualism (which was to be the swan song of opera in any case). This particular ensemble (and this soprano) specialize in `authentic' baroque performances and is most excellent:

\url{https://www.youtube.com/watch?v=H2Ks_TurCUo}


This song I first heard on one of the finest evenings of my life:

\url{https://www.youtube.com/watch?v=Bofbxcz4bqw}

\hfill

\end{sffamily}
\end{footnotesize}

